\chapter*{Résumé}
\addcontentsline{toc}{chapter}{Résumé}

L'adoption du cloud computing dans les institutions académiques tunisiennes est freinée
par l'absence de portails en libre-service pour l'infrastructure nationale
Huawei Cloud Stack (HCS) gérée par le Ministère de l'Enseignement Supérieur et de la
Recherche Scientifique (MESRS). Les chercheurs et étudiants souhaitant accéder à des
machines virtuelles à la demande doivent recourir à des processus administratifs
manuels et chronophages, sans visibilité sur l'avancement du provisionnement.

Ce projet présente \textbf{VM Marketplace}, une plateforme cloud full-stack en
libre-service, conçue et développée lors d'un stage de six mois chez Huawei
Technologies Tunisie. La plateforme permet aux utilisateurs authentifiés de
configurer, payer et provisionner des machines virtuelles et des clusters Kubernetes
directement sur l'infrastructure HCS via une interface web intuitive. Ses
fonctionnalités clés comprennent : un assistant de configuration de VM en plusieurs
étapes avec suivi en temps réel via Server-Sent Events, un moteur de recommandation
IA basé sur le modèle de langage Llama~3.1, un terminal SSH dans le navigateur,
une installation automatisée de Netdata pour la supervision, un gestionnaire complet
du cycle de vie des clusters k3s, un contrôleur d'autoguérison hub-spoke et des outils
de synchronisation entre la base de données et l'état Terraform cloud.

L'architecture repose sur trois services indépendants : un frontend Next.js, un backend
FastAPI et un moteur d'inférence IA distinct. Le provisionnement d'infrastructure est
géré par Terraform avec le fournisseur \texttt{huaweicloud/hcs}. Les clusters
Kubernetes sont initialisés avec k3s (distribution légère de Rancher Labs). Un schéma
réseau de type \og{}maître-comme-passerelle-NAT\fg{} --- combinant
\texttt{source\_dest\_check=false} au niveau de l'hyperviseur avec le masquage IP Linux
--- a été conçu et validé afin de fournir aux nœuds workers un accès internet malgré
les contraintes de quota d'EIP de l'environnement HCS partagé.

La validation sur l'infrastructure HCS de production a confirmé le provisionnement
complet de machines virtuelles et d'un cluster à deux nœuds sous k3s v1.31.4+k3s1,
avec les deux nœuds en statut \texttt{Ready}.

\bigskip
\noindent\textbf{Mots-clés :} Cloud computing, Huawei Cloud Stack, Machines virtuelles,
Kubernetes, k3s, Infrastructure as Code, Terraform, Recommandation IA, FastAPI, Next.js,
Portail libre-service.

\clearpage
